%==========================================================================
% Chapter: Methodology
% Requires in the preamble of bachelor.tex:  \usepackage{amsmath}
% Figures marked with \placeholderfig are to be replaced on Overleaf by
%   \includefig{<width>}{<file name>}{<caption>}{<label>}
% once the corresponding image has been placed in the images/ folder.
%==========================================================================

% Temporary figure placeholder: same argument order as \includefig plus a
% short note describing what the image must show.
\providecommand{\placeholderfig}[5]{
    \begin{figure}[ht]
     \centering
      \fbox{\parbox[c][0.28\textheight][c]{#1\textwidth}{\centering
        \textbf{Figure placeholder}\\[2mm]
        \texttt{\detokenize{images/#2}}\\[2mm]
        \small #5}}
      \caption{#3}
      \label{#4}
    \end{figure}
}

\chapter{Model Predictive Control Design for the Inverted Pendulum on a Cart}
\label{chap:methodology}
\chaptermark{MPC Design}

This chapter describes how the \ac{mpc} stabilising
controller for the inverted pendulum on a cart was designed, implemented, and
prepared for experimental evaluation. Section~\ref{sec:requirements} states the
control problem and the requirements derived from the hardware.
Section~\ref{sec:platform} describes the experimental platform, and
Section~\ref{sec:model} derives the nonlinear and linearised models of the
plant. Section~\ref{sec:baseline} presents the existing control structure of
the laboratory course into which the new controller is integrated, and
Section~\ref{sec:selection} justifies the choice of a linear, constrained
\ac{mpc}. The controller itself is formulated in Section~\ref{sec:mpc}, its
implementation in Simulink is described in Section~\ref{sec:implementation},
and the procedure used to evaluate it is defined in
Section~\ref{sec:evaluation}.

The controller has so far been designed and verified in simulation only. Its
execution on the laboratory hardware and its comparison with the \ac{lqr} and \ac{pid} controllers are part of the planned work described in
Section~\ref{sec:evaluation}.

%--------------------------------------------------------------------------
\section{Problem Statement and Design Requirements}
\label{sec:requirements}
%--------------------------------------------------------------------------

The plant is a pendulum rod that pivots freely on a cart driven along a linear
axis. Only the cart is actuated; the pendulum angle can be influenced solely
through the motion of the cart, so the system is underactuated. The upright
position of the pendulum is an unstable equilibrium.

The control task considered in this thesis is the stabilising (balancing)
phase: once the pendulum has been brought close to the upright position, the
controller must
\begin{enumerate}
  \item keep the pendulum upright, $\varphi \rightarrow 0$, in the presence of
        initial deviations and disturbances;
  \item move the cart to a commanded position set-point $r$ without steady-state
        offset;
  \item respect the physical limits of the drive and the axis while doing so.
\end{enumerate}
The swing-up of the pendulum from the hanging position is performed by the
existing energy-based controller of the laboratory course
(Section~\ref{sec:baseline}) and is not redesigned here.

The requirements in Table~\ref{tab:requirements} follow from the hardware
specifications and from the existing software environment. The limits on cart
speed and acceleration are those of the internal drive control of the
CO3620-2G~\cite{ln_manual}. The travel limit and the update rate result from
the design choices described in Sections~\ref{subsec:constraints}
and~\ref{subsec:horizons}.

\begin{table}[ht]
\centering
\caption{Design requirements for the stabilising controller.}
\label{tab:requirements}
\begin{tabular}{p{0.05\textwidth}p{0.5\textwidth}p{0.35\textwidth}}
\toprule
No. & Requirement & Origin \\
\midrule
R1 & Stabilise the pendulum about the upright equilibrium. & Control objective \\
R2 & Track a cart position set-point without steady-state offset. & Control objective \\
R3 & Cart acceleration command limited to $|a| \le 10~\mathrm{m/s^2}$. & Drive specification~\cite{ln_manual} \\
R4 & Cart speed limited to $|\dot{x}| \le 0.5~\mathrm{m/s}$. & Drive specification~\cite{ln_manual} \\
R5 & Cart travel limited to $|x| \le 0.15~\mathrm{m}$ about the centre of the axis. & Design choice within the 400~mm traversing distance \\
R6 & Controller executable on the CO3620-2A digital controller within the time budget of its control routine. & Target hardware \\
R7 & No modification of the existing hardware blocks, swing-up controller, mode logic, or state observer of the course model. & Integration constraint \\
\bottomrule
\end{tabular}
\end{table}

%--------------------------------------------------------------------------
\section{Experimental Platform}
\label{sec:platform}
%--------------------------------------------------------------------------

The work is based on the Lucas-Nülle training system for automatic control of a
position and angle control system with MATLAB/Simulink (course
SO2802-3H)~\cite{ln_manual}. The system consists of three units: the position
and angle control system with inverted pendulum CO3620-2G, the digital
universal controller CO3620-2A, and the MATLAB programming interface and
toolbox CO3636-3D. Figure~\ref{fig:setup} shows the setup.

\placeholderfig{0.8}{ln_setup_photo.png}
  {Laboratory setup: CO3620-2G position and angle control system with inverted pendulum and CO3620-2A digital universal controller.}
  {fig:setup}
  {Photograph of the complete lab setup (cart, linear axis, pendulum rod, CO3620-2A controller unit). Please provide.}

\subsection{Position and Angle Control System CO3620-2G}
\label{subsec:co3620-2g}

The CO3620-2G consists of a linear axis with belt drive and a linear guidance
system with recirculating ball bearings. The axis offers a traversing distance
of up to 400~mm. The cart is driven by a hybrid stepping motor with a power of
up to 100~W and integrated speed and current control, which limits the cart
speed to 0.5~m/s and the acceleration to 10~m/s$^2$. The set-point speed of the
drive is commanded through an analogue signal of $\pm 10$~V. A pendulum rod is
mounted on a pivot on the cart~\cite{ln_manual}.

The cart position is measured with an incremental encoder on the motor shaft.
The linear axis has a feed of 54~mm per motor revolution, and the
\emph{Encoder Calculation} block of the Lucas-Nülle toolbox delivers 16000
counts per revolution, which gives a position resolution of
$0.054/16000~\mathrm{m} = 3.375~\mu\mathrm{m}$. The pendulum angle is measured
by an encoder read over the \acf{can} bus with a
resolution of 65535 counts per revolution~\cite{ln_manual}.

\subsection{Digital Universal Controller and Toolchain}
\label{subsec:co3620-2a}

The control algorithm runs on the digital universal controller CO3620-2A,
which is programmed through the CO3636-3D interface from Simulink. The
Simulink model is translated into C code with Embedded Coder, compiled, and
downloaded to the controller; a USB connection exchanges parameters and
measurement data with MATLAB during operation~\cite{ln_manual}. The programming
adapter is the C2000 F28335 XDS100V2, and the controller is therefore based on
a Texas Instruments TMS320F28335 \acf{dsp}.

In the course models, the \emph{LucasNuelleHardware} block executes a
subsystem called \emph{Control-Routine} as an \acf{isr} at a frequency of $f_{\mathrm{ISR}} = 8~\mathrm{kHz}$, i.e.\ with a
period of $T_{\mathrm{ISR}} = 125~\mu\mathrm{s}$. All blocks inside the
Control-Routine inherit this sample time.

\subsection{Signal Interface}
\label{subsec:signals}

Table~\ref{tab:signals} lists the signals that connect the Control-Routine to
the hardware. The angle reading is reset with the pendulum hanging, after
which the upright position corresponds to $\varphi = 0$. The speed set-point
$y_1$ is scaled such that 10~V corresponds to 0.5~m/s, i.e.\ a factor of
20~V/(m/s).

\begin{table}[ht]
\centering
\caption{Hardware signals used by the controller.}
\label{tab:signals}
\begin{tabular}{llll}
\toprule
Signal & Direction & Unit & Scaling / remark \\
\midrule
Cart position $x$ & Input & m & $0.054/16000$~m per count \\
Pendulum angle $\varphi$ & Input & rad & $2\pi/65535$~rad per count; upright $= 0$ \\
Speed set-point $y_1$ & Output & V & $\pm 10$~V $\hat{=}$ $\pm 0.5$~m/s \\
Drive enable $Q_2$ & Parameter & -- & Enables the drive \\
Position set-point & Parameter & mm & Set\_Pos, set from the PC \\
\bottomrule
\end{tabular}
\end{table}

%--------------------------------------------------------------------------
\section{Mathematical Model}
\label{sec:model}
%--------------------------------------------------------------------------

\subsection{Nonlinear Equations of Motion}
\label{subsec:nonlinear}

The model follows the derivation in the course documentation~\cite{ln_manual}.
The notation is summarised in Table~\ref{tab:params}. The angle $\theta$ is
measured from the hanging position, so that $\theta = \pi$ corresponds to the
upright pendulum. Figure~\ref{fig:fbd} shows the forces acting on the cart and
the pendulum.

\placeholderfig{0.6}{cart_pendulum_fbd.png}
  {Free-body diagram of the cart and the pendulum, adapted from~\cite{ln_manual}.}
  {fig:fbd}
  {Free-body diagram: cart (mass $M$, force $F$, friction), pendulum (mass $m$, distance $l$, angle $\theta$), reaction forces $N$ and $P$, gravity. Please redraw (e.g.\ in draw.io/TikZ) rather than copying the manual figure.}

The horizontal force balance of the cart, with viscous friction coefficient $b$,
is
\begin{equation}
  M\ddot{x} = F - N - b\dot{x},
  \label{eq:cart}
\end{equation}
and the torque balance of the pendulum about its centre of gravity is
\begin{equation}
  I\ddot{\theta} = -Nl\cos\theta - Pl\sin\theta,
  \label{eq:pend_torque}
\end{equation}
where $N$ and $P$ are the horizontal and vertical reaction forces at the pivot.
With the position of the centre of gravity
$x_p = x + l\sin\theta$, $y_p = -l\cos\theta$, the force balances of the
pendulum, $N = m\ddot{x}_p$ and $P - mg = m\ddot{y}_p$, yield
\begin{align}
  N &= m\ddot{x} + ml\ddot{\theta}\cos\theta - ml\dot{\theta}^2\sin\theta,
  \label{eq:N}\\
  P &= ml\dot{\theta}^2\cos\theta + ml\ddot{\theta}\sin\theta + mg.
  \label{eq:P}
\end{align}
Substituting \eqref{eq:N} and \eqref{eq:P} into \eqref{eq:pend_torque}, the
centripetal terms cancel and the pendulum dynamics become
\begin{equation}
  \left(I + ml^2\right)\ddot{\theta} = -ml\,\ddot{x}\cos\theta - mgl\sin\theta .
  \label{eq:pend_nl}
\end{equation}

\subsection{Cart Acceleration as Control Input}
\label{subsec:input}

The drive of the CO3620-2G contains its own closed speed loop: the analogue
signal $y_1$ commands a cart speed, not a force. The force $F$ in
\eqref{eq:cart}, the cart mass $M$, and the friction $b$ are therefore not
accessible to the controller. Following the course model~\cite{ln_manual}, the
cart acceleration $a = \ddot{x}$ is taken as the control input of the model.
Equation~\eqref{eq:cart} is then no longer needed, and the cart reduces to a
double integrator, $\ddot{x} = a$. The controller output is converted into a
speed set-point by integration,
\begin{equation}
  v_{\mathrm{set}}(t) = \int_{t_0}^{t} a(\tau)\,\mathrm{d}\tau, \qquad
  y_1 = 20~\tfrac{\mathrm{V}}{\mathrm{m/s}} \cdot v_{\mathrm{set}},
  \label{eq:integrator}
\end{equation}
where $t_0$ is the instant at which the balancing controller is activated. This
structure assumes that the drive tracks the speed set-point closely; the
deviation between the commanded and the actual acceleration is one of the
sources of model mismatch discussed in Section~\ref{subsec:simmodel}.

\subsection{Linearisation about the Upright Equilibrium}
\label{subsec:linear}

The angle $\varphi = \theta - \pi$ is measured from the upright position, so
that $\sin\theta = -\sin\varphi$ and $\cos\theta = -\cos\varphi$. For small
deviations, $\sin\varphi \approx \varphi$ and $\cos\varphi \approx 1$, and
\eqref{eq:pend_nl} becomes~\cite{ln_manual}
\begin{equation}
  \left(I + ml^2\right)\ddot{\varphi} - mgl\,\varphi = ml\,a .
  \label{eq:pend_lin}
\end{equation}
With the state vector $\mathbf{x} = [x,\ \dot{x},\ \varphi,\ \dot{\varphi}]^{\mathsf{T}}$,
the input $u = a$, and the measured outputs $x$ and $\varphi$, the
continuous-time state-space model is
\begin{equation}
  \dot{\mathbf{x}} = A\mathbf{x} + Bu, \qquad \mathbf{y} = C\mathbf{x},
  \label{eq:ss}
\end{equation}
\begin{equation}
  A = \begin{bmatrix}
        0 & 1 & 0 & 0\\
        0 & 0 & 0 & 0\\
        0 & 0 & 0 & 1\\
        0 & 0 & \dfrac{mgl}{I+ml^2} & 0
      \end{bmatrix}, \quad
  B = \begin{bmatrix} 0\\ 1\\ 0\\ \dfrac{ml}{I+ml^2} \end{bmatrix}, \quad
  C = \begin{bmatrix} 1 & 0 & 0 & 0\\ 0 & 0 & 1 & 0 \end{bmatrix}.
  \label{eq:ABC}
\end{equation}

\subsection{Model Parameters}
\label{subsec:params}

The parameters used in the course initialisation script are listed in
Table~\ref{tab:params}. The moment of inertia $I = ml^2/3$ corresponds to a
slender rod of length $2l$ about its centre of gravity.

\begin{table}[ht]
\centering
\caption{Model parameters of the pendulum (course initialisation script).}
\label{tab:params}
\begin{tabular}{llll}
\toprule
Symbol & Description & Value & Unit \\
\midrule
$m$ & Mass of the pendulum & 0.3 & kg \\
$l$ & Distance from pivot to centre of gravity & 0.165 & m \\
$I$ & Moment of inertia about the centre of gravity, $ml^2/3$ & $2.72\times10^{-3}$ & kg\,m$^2$ \\
$g$ & Gravitational acceleration & 9.81 & m/s$^2$ \\
\bottomrule
\end{tabular}
\end{table}

With these values,
\begin{equation}
  \frac{mgl}{I+ml^2} = \frac{3g}{4l} = 44.59~\mathrm{s^{-2}}, \qquad
  \frac{ml}{I+ml^2} = \frac{3}{4l} = 4.545~\mathrm{m^{-1}} .
  \label{eq:numeric}
\end{equation}
The eigenvalues of $A$ are $\pm 6.678~\mathrm{rad/s}$ and a double eigenvalue at
zero. The positive eigenvalue confirms the instability of the upright
equilibrium; its time constant is $1/6.678 \approx 0.15$~s. The double zero
eigenvalue represents the cart, which behaves as a double integrator. The
controllability matrix of $(A,B)$ and the observability matrix of $(A,C)$ both
have full rank 4, so the system is controllable from the cart acceleration and
observable from the measurements of $x$ and $\varphi$.

\subsection{Discretisation}
\label{subsec:discrete}

The \ac{mpc} operates on a discrete-time model. The model \eqref{eq:ss} was
discretised with a \acf{zoh} at the controller sample time
$T_s = 10$~ms (Section~\ref{subsec:horizons}), which gives
$\mathbf{x}(k+1) = A_d\mathbf{x}(k) + B_d u(k)$ with
\begin{equation}
  A_d = \begin{bmatrix}
    1 & 0.01 & 0 & 0\\
    0 & 1 & 0 & 0\\
    0 & 0 & 1.0022 & 0.0100\\
    0 & 0 & 0.4462 & 1.0022
  \end{bmatrix}, \qquad
  B_d = \begin{bmatrix}
    5.0\times10^{-5}\\ 0.01\\ 2.27\times10^{-4}\\ 0.0455
  \end{bmatrix}.
  \label{eq:AdBd}
\end{equation}
The eigenvalues of $A_d$ are $1.0691$, $0.9354$, and a double eigenvalue at
1; the first lies outside the unit circle and corresponds to the unstable
continuous-time eigenvalue, $e^{6.678\,T_s} = 1.0691$.

%--------------------------------------------------------------------------
\section{Existing Control Structure of the Course Model}
\label{sec:baseline}
%--------------------------------------------------------------------------

The \ac{mpc} is integrated into the course model \emph{A62\_SwingUp\_Control},
which brings the pendulum from the hanging to the upright position and then
balances it~\cite{ln_manual}. Its Control-Routine contains three controllers,
a state observer, and a switching logic (Figure~\ref{fig:a62}). Only the
balancing controller is replaced in this work (requirement R7); the other
elements are described here because they define the conditions under which
the \ac{mpc} operates.

\placeholderfig{0.95}{a62_control_routine.png}
  {Structure of the Control-Routine of the course model with the swing-up controller, the position controller, the state space controller, and the controller selection logic.}
  {fig:a62}
  {Screenshot of the Simulink subsystem \emph{Control-Routine} (from A62\_MPC\_Sim.slx or A62\_SwingUp\_Control.slx), showing the three controllers, the two Controller Selection blocks and the two Switch blocks. Please provide.}

\subsection{Swing-Up Controller}
\label{subsec:swingup}

The swing-up controller is based on the energy of the pendulum. With the angle
measured from the upright position, the energy
\begin{equation}
  E = \tfrac{1}{2}I\dot{\varphi}^2 + mgl\left(\cos\varphi - 1\right)
  \label{eq:energy}
\end{equation}
is zero when the pendulum is at rest in the upright position. The course
derives from a Lyapunov argument that energy is pumped into the pendulum when
the cart is accelerated according to the sign of $E\,\dot{\varphi}\cos\varphi$
~\cite{ln_manual}. In the implementation, this law is realised as a
bang-bang position set-point: the cart is commanded to $+150$~mm if
$-\operatorname{sign}(E\,\dot{\varphi}\cos\varphi) > 0$ and to $-150$~mm
otherwise. The angular velocity is obtained with a filtered derivative
(derivative gain 1, filter coefficient $N = 100$).

\subsection{Position Controller}
\label{subsec:poscontrol}

The position controller computes the error $e = r - x$, multiplies it by a
gain of 4.17~s$^{-1}$, passes the result through a first-order lag with time
constant 0.17~s, and scales it by 20~V/(m/s) to obtain the speed set-point
$y_1$. It follows the bang-bang set-point of the swing-up controller and is
also the fallback when the pendulum leaves the region in which balancing is
possible.

\subsection{State Space Controller and Observer}
\label{subsec:lqr}

The balancing controller of the course is a state-feedback controller with a
static prefilter,
\begin{equation}
  a = \bar{N}r - K\hat{\mathbf{x}},
  \label{eq:lqr}
\end{equation}
followed by the integrator \eqref{eq:integrator}. The gain $K$ is the
continuous-time \ac{lqr} gain for the model \eqref{eq:ss} with
$Q = \operatorname{diag}(1000, 0, 1000, 0)$ and $R = 1$, and $\bar{N}$ is chosen
so that the cart position has unit static gain. The resulting values are
\begin{equation}
  K = \begin{bmatrix} -31.62 & -20.14 & 72.76 & 9.81 \end{bmatrix}, \qquad
  \bar{N} = -31.62,
  \label{eq:Kvalues}
\end{equation}
with closed-loop eigenvalues at $-9.88 \pm 7.41j$ and $-2.35 \pm 1.94j$. The
negative position gain reflects the non-minimum-phase behaviour of the
position response: to move the cart in one direction, the controller first
moves it in the opposite direction so that the pendulum tilts towards the
target. This controller serves as the \ac{lqr} reference in the planned
comparison (Section~\ref{subsec:benchmark}).

Only $x$ and $\varphi$ are measured, so the full state is reconstructed by a
Luenberger observer. Because the drive does not reproduce the commanded
acceleration exactly, the observer first models the drive: the speed
set-point, limited to $\pm 0.5$~m/s, is compared with the modelled cart speed,
and the difference is fed to a proportional--integral controller ($K_p = 100$,
$K_i = 30$, output limited to $\pm 10~\mathrm{m/s^2}$) whose output
$a_{\mathrm{model}}$ is the modelled cart acceleration. The observer then
evaluates
\begin{equation}
  \dot{\hat{\mathbf{x}}} = A\hat{\mathbf{x}} + B a_{\mathrm{model}}
    + L\left(\mathbf{y} - C\hat{\mathbf{x}}\right),
  \label{eq:observer}
\end{equation}
with the gain $L$ obtained by pole placement for the eigenvalues
\begin{equation}
  \lambda_{\mathrm{obs}} = -\frac{1}{T_L}\left\{1,\ \frac{1}{1.01},\ \frac{1}{1.02},\ \frac{1}{1.03}\right\},
  \qquad T_L = 5~\mathrm{ms},
  \label{eq:obspoles}
\end{equation}
i.e.\ between $-194$ and $-200~\mathrm{rad/s}$. The integrations are carried out with the
forward Euler method at the rate of the Control-Routine.

\subsection{Controller Selection Logic}
\label{subsec:modelogic}

Two latching logic blocks select the active controller
(Table~\ref{tab:modelogic}). The first records that the pendulum has reached
the top; the second activates the balancing controller only within a capture
region and keeps it active within a larger holding region. Outside the holding
region, the position controller takes over.

\begin{table}[ht]
\centering
\caption{Controller selection conditions of the course model.}
\label{tab:modelogic}
\begin{tabular}{p{0.3\textwidth}p{0.6\textwidth}}
\toprule
Condition & Definition \\
\midrule
Upright reached (latched) & $|\varphi| < 2^\circ$; cleared on a rising edge of $Q_2$ \\
Capture (switch to balancing) & $|\varphi| < 15^\circ$, $|\dot{\varphi}| < 130^\circ/\mathrm{s}$, and $|x| < 30$~mm \\
Hold (stay in balancing) & $|\varphi| < 20^\circ$ and $|x| < 150$~mm \\
\bottomrule
\end{tabular}
\end{table}

%--------------------------------------------------------------------------
\section{Selection of the Control Approach}
\label{sec:selection}
%--------------------------------------------------------------------------

The course already contains a working \ac{lqr}. The \ac{lqr} is optimal for
the unconstrained linear problem, but the limits R3--R5 in
Table~\ref{tab:requirements} are not part of its design. In the course model,
the cart speed and acceleration are limited only downstream of the controller,
by saturation in the drive, and the travel limit is enforced only by the
controller selection logic, which abandons balancing when $|x| \ge 150$~mm.
A \ac{pid} controller acting on the pendulum angle alone cannot regulate the
cart position at the same time; the course therefore uses cascaded loops for
this task~\cite{ln_manual}, and such loops again treat the limits only through
saturation.

\ac{mpc} addresses these points directly. The limits enter the optimisation
as constraints, so that the predicted trajectory remains admissible over the
prediction horizon instead of being clipped afterwards~\cite{rawlings2000}.
Cart and pendulum are handled in one multivariable problem. For a linear
model, a quadratic cost and linear constraints, the optimisation is a convex \ac{qp}~\cite{rawlings2000}. Moreover, when the terminal
cost is the solution of the \acf{dare} and no constraint is active, the \ac{mpc} control law coincides
with the infinite-horizon \ac{lqr}~\cite{rawlings2000}. This property is
used here to start the tuning from a controller equivalent to the
\ac{lqr}, so that the effect of the constraints can later be separated from
the effect of the weights. This addresses the limitation of the comparison
in~\cite{minouchehr2015}, where an unconstrained \ac{mpc} was compared with an
\ac{lqr}.

Two alternative \ac{mpc} variants were considered and not pursued.
A nonlinear \ac{mpc} based on \eqref{eq:pend_nl} would describe large angles
more accurately, but the controller only operates within the holding region of
Table~\ref{tab:modelogic}, and a nonlinear program is considerably more
expensive to solve on the target \ac{dsp}~\cite{rawlings2000}. Explicit
\ac{mpc}, which precomputes the control law offline, is practical only for
small problems with few constraints and short horizons~\cite{wang2010}; the
horizon used here (Section~\ref{subsec:horizons}) exceeds this range. A linear
\ac{mpc} with an online \ac{qp} solver was therefore selected.

%--------------------------------------------------------------------------
\section{MPC Formulation}
\label{sec:mpc}
%--------------------------------------------------------------------------

The controller was designed with the Model Predictive Control Toolbox of
MATLAB, starting from the MathWorks example on the control of an inverted
pendulum on a cart~\cite{mathworks_pendcart}, and adapted to the model and
limits of the CO3620-2G.

\subsection{Prediction Model and Signals}
\label{subsec:predmodel}

The prediction model is the discrete-time model \eqref{eq:AdBd}. Its single
manipulated variable is the cart acceleration $u = a$. The cart position $x$
and the pendulum angle $\varphi$ are declared as measured outputs, and the cart
speed $\dot{x}$ is added as an unmeasured output so that a constraint can be
imposed on it. To implement the terminal cost (Section~\ref{subsec:cost}),
four further unmeasured outputs $\mathbf{z} = L_c\mathbf{x}$ are appended,
where $L_c$ is the Cholesky factor of the terminal weight,
$L_c^{\mathsf{T}}L_c = P$. The output vector of the prediction model is thus
\begin{equation}
  \mathbf{y}_{\mathrm{MPC}} =
  \begin{bmatrix} x & \varphi & \dot{x} & \mathbf{z}^{\mathsf{T}} \end{bmatrix}^{\mathsf{T}}
  = \begin{bmatrix} C_m \\ L_c \end{bmatrix}\mathbf{x}, \qquad
  C_m = \begin{bmatrix} 1&0&0&0\\ 0&0&1&0\\ 0&1&0&0 \end{bmatrix}.
  \label{eq:mpcoutputs}
\end{equation}

\subsection{Cost Function}
\label{subsec:cost}

At each sample $k$, the controller minimises
\begin{equation}
  J(k) = \sum_{i=1}^{N_p-1} \tilde{\mathbf{x}}(k+i|k)^{\mathsf{T}} Q_d\,
         \tilde{\mathbf{x}}(k+i|k)
       + \tilde{\mathbf{x}}(k+N_p|k)^{\mathsf{T}} P\,
         \tilde{\mathbf{x}}(k+N_p|k)
       + \sum_{i=0}^{N_p-1} R_d\, u(k+i|k)^2
       + \rho\,\varepsilon^2,
  \label{eq:cost}
\end{equation}
where $\tilde{\mathbf{x}} = \mathbf{x} - \mathbf{x}_{\mathrm{ref}}$ is the
deviation from the reference state $\mathbf{x}_{\mathrm{ref}} = [r,\ 0,\ 0,\ 0]^{\mathsf{T}}$,
$N_p$ is the prediction horizon, $\varepsilon \ge 0$ is the slack variable of
the soft constraints (Section~\ref{subsec:constraints}), and $\rho = 10^5$ is
its penalty. The input is free for the first $N_c$ steps of the horizon and
held constant afterwards.

The stage weights are
\begin{equation}
  Q_d = \tfrac{1}{90}\operatorname{diag}(q_x,\ q_v,\ q_\varphi,\ 0), \qquad
  R_d = \tfrac{1}{90},
  \label{eq:QdRd}
\end{equation}
where the common factor $1/90$ rescales the cost without changing its
minimiser. The terminal weight $P$ is the solution of the \ac{dare} for
$(A_d, B_d, Q_d, R_d)$, so that the terminal term approximates the cost of the
infinite horizon beyond step $N_p$. In the toolbox, the stage cost is realised
through weights on the outputs $x$, $\varphi$, and $\dot{x}$ for steps $1$ to
$N_p - 1$, and the terminal cost through unit weights on $\mathbf{z}$ at step
$N_p$ only, since $\tilde{\mathbf{z}}^{\mathsf{T}}\tilde{\mathbf{z}} =
\tilde{\mathbf{x}}^{\mathsf{T}}P\,\tilde{\mathbf{x}}$. The reference for
$\mathbf{z}$ is accordingly $L_c\mathbf{x}_{\mathrm{ref}}$; without it, the
terminal term would pull the cart towards $x = 0$ and cause a steady-state
offset. The weights are specified in the toolbox as weights on scaled
variables (scale factors 0.3~m for $x$, 0.7~rad for $\varphi$, 1~m/s for
$\dot{x}$, and 20~m/s$^2$ for $u$) and were chosen such that the cost in
physical units equals \eqref{eq:cost}. No weight is placed on the rate of
change of the input.

\subsection{Constraints}
\label{subsec:constraints}

The optimisation is subject to the prediction model and to
\begin{align}
  -10~\mathrm{m/s^2} \le\ & u(k+i|k) \le 10~\mathrm{m/s^2},
  \label{eq:con_u}\\
  -0.15~\mathrm{m} - \varepsilon \le\ & x(k+i|k) \le 0.15~\mathrm{m} + \varepsilon,
  \label{eq:con_x}\\
  -0.5~\mathrm{m/s} - \varepsilon \le\ & \dot{x}(k+i|k) \le 0.5~\mathrm{m/s} + \varepsilon,
  \label{eq:con_v}
\end{align}
with the slack expressed in scaled units (equal concern for relaxation of 1
for both outputs). The acceleration limit \eqref{eq:con_u} corresponds to the
drive specification and is a hard constraint. The speed and travel limits are
soft constraints: they are respected whenever possible, but they can be
violated at a high cost if a disturbance would otherwise make the \ac{qp}
infeasible~\cite{rawlings2000}. The travel limit of $\pm 0.15$~m equals the
holding limit of the controller selection logic (Table~\ref{tab:modelogic});
a predicted trajectory that would cause the logic to abandon balancing is
therefore penalised by the controller itself. The pendulum angle is weighted but not
constrained.

\subsection{Sample Time and Horizons}
\label{subsec:horizons}

The \ac{mpc} cannot solve a \ac{qp} at the 8~kHz rate of the Control-Routine.
A controller sample time of $T_s = 10$~ms was therefore chosen, so that the
\ac{mpc} is evaluated once every $n_{\mathrm{MPC}} = T_s f_{\mathrm{ISR}} = 80$
calls of the Control-Routine. This sample time is about 15 times shorter than
the time constant of the unstable mode (0.15~s, Section~\ref{subsec:params}).
The prediction horizon is $N_p = 50$ steps, i.e.\ 0.5~s, which covers about
3.3 time constants of the unstable mode. The control horizon is $N_c = 5$,
which limits the \ac{qp} to five decision variables plus the slack variable.

\subsection{Weight Tuning}
\label{subsec:tuning}

The tuning started from weights equivalent to the course \ac{lqr},
$Q = \operatorname{diag}(1000, 0, 1000, 0)$ and $R = 1$, and the position and
speed weights were then increased. Four candidates were simulated for a
100~mm position step with the linear model and the constraints active
(Table~\ref{tab:candidates}).

% NOTE: The values in this table were obtained by running Calculations.m
% (loop over the four candidates, 300 steps = 3 s, stepinfo with 2 %
% settling threshold and 10-90 % rise time). Overshoot and undershoot are
% given in % of the 100 mm step (numerically equal to mm here).
% Please confirm that the selection argument below matches your reasoning.
\begin{table}[ht]
\centering
\caption{Weight candidates and their simulated response to a 100~mm position step (linear prediction model, $T_s = 10$~ms, $N_p = 50$, $N_c = 5$, $R = 1$).}
\label{tab:candidates}
\small
\setlength{\tabcolsep}{4pt}
\begin{tabular}{lccccccc}
\toprule
Candidate & $q_x$ & $q_v$ & $q_\varphi$ & Settling & Overshoot & Undershoot & $\max|a|$ \\
          &       &       &             & (2\,\%) [s] & [\%] & [\%] & [m/s$^2$] \\
\midrule
LQR-equivalent & 1000 & 0 & 1000 & 2.42 & 3.6 & 5.6 & 2.77 \\
A & 3000 & 0 & 1000 & 1.93 & 3.3 & 8.1 & 4.33 \\
B & 10000 & 0 & 1000 & 1.46 & 2.6 & 12.1 & 7.20 \\
D & 5000 & 100 & 1000 & 1.16 & 1.9 & 9.1 & 5.24 \\
\bottomrule
\end{tabular}
\end{table}

Candidate D gives the shortest settling time and the smallest overshoot of the
four candidates. Its larger undershoot and peak acceleration remain within the
travel and acceleration limits, and it was retained as the final weighting,
$q_x = 5000$, $q_v = 100$, $q_\varphi = 1000$. The step responses of the
candidates are shown in Figure~\ref{fig:candidates}; the detailed evaluation of
the final controller is presented in the Results chapter.

\placeholderfig{0.85}{mpc_candidates_step.png}
  {Simulated response of the four weight candidates and the course \acs{lqr} to a 100~mm position step: cart position, pendulum angle, and cart acceleration.}
  {fig:candidates}
  {The 3-panel figure produced by the candidate loop in Calculations.m (x [mm], $\varphi$ [deg], a [m/s$^2$] vs.\ t; four candidates + course LQR dashed). Please export it as PNG and provide.}

\subsection{State Estimation}
\label{subsec:estimation}

The toolbox provides a built-in state estimator with an output disturbance
model. In this design, it is disabled: the output disturbance model is
removed, and the controller is configured for custom state estimation. At each
\ac{mpc} sample, the plant state of the controller is set to the estimate
$\hat{\mathbf{x}}$ of the existing Luenberger observer \eqref{eq:observer}.
This keeps the observer of the course unchanged (requirement R7) and gives the
\ac{mpc} and the \ac{lqr} the same state information, so that differences in
performance can be attributed to the control law. Because no disturbance model
is included, the controller does not contain integral action; offset-free
position tracking relies on the accuracy of the model and of the estimate.

\subsection{Solution of the Optimisation Problem}
\label{subsec:qp}

The cost \eqref{eq:cost}, the linear model, and the linear constraints
\eqref{eq:con_u}--\eqref{eq:con_v} form a convex \ac{qp}, which is solved at
every \ac{mpc} sample with the active-set solver of the Model Predictive
Control Toolbox. Only the first element $u(k|k)$ of the optimal input sequence
is applied, and the problem is solved again at the next sample with a new
state estimate (receding horizon~\cite{rawlings2000}).

%--------------------------------------------------------------------------
\section{Implementation}
\label{sec:implementation}
%--------------------------------------------------------------------------

\subsection{Integration into the Control-Routine}
\label{subsec:integration}

The \ac{mpc} replaces the prefilter, the state feedback, and the summation of
\eqref{eq:lqr} inside the \emph{State Space Controller} subsystem; the
integrator \eqref{eq:integrator}, the observer, and all hardware blocks are
unchanged (Figure~\ref{fig:sscontroller}). A constant \texttt{useMPC} and a
switch select either the \ac{mpc} or the original \ac{lqr} output as the
input of the integrator. This allows both controllers to be run in the same
model and returns the system to the course \ac{lqr} by changing a single
parameter.

\placeholderfig{0.9}{state_space_controller_mpc.png}
  {Modified State Space Controller subsystem: the \acs{mpc} block and the original \acs{lqr} path, selected by the \texttt{useMPC} switch, both feeding the acceleration-to-speed integrator.}
  {fig:sscontroller}
  {Screenshot of Control-Routine $\rightarrow$ State Space Controller in A62\_MPC\_Sim.slx, showing the observer, the block ``MPC (candidate D)'', K, Nbar, Sum, the ``LQR or MPC'' switch, useMPC, the Discrete-Time Integrator and the gain 20. Please provide.}

\begin{sloppypar}
The controller is implemented in a MATLAB Function block that calls
\texttt{mpcmoveCodeGeneration}, the code-generation form of the toolbox solver.
The controller object is converted offline into the constant structures
required by this function with \texttt{getCodeGenerationData}. The block has
three inputs: the observer estimate $\hat{\mathbf{x}}$, the position
set-point, and the enable signal of the balancing mode. Its operation is
summarised in Algorithm~\ref{alg:mpcblock}.
\end{sloppypar}

\begin{algorithm}[ht]
\caption{Execution of the \acs{mpc} block in each call of the Control-Routine.}
\label{alg:mpcblock}
\begin{algorithmic}[1]
\IF{balancing is enabled}
  \IF{$c = 0$}
    \STATE set the controller state to $\hat{\mathbf{x}}$ (custom estimation)
    \STATE set the references $[r,\ 0,\ 0,\ (L_c\mathbf{x}_{\mathrm{ref}})^{\mathsf{T}}]$ with $r$ in metres
    \STATE solve the \ac{qp} and store the first input $a_{\mathrm{hold}} \leftarrow u(k|k)$
  \ENDIF
  \STATE $c \leftarrow c + 1$; \textbf{if} $c \ge n_{\mathrm{MPC}}$ \textbf{then} $c \leftarrow 0$
\ELSE
  \STATE $c \leftarrow 0$, $a_{\mathrm{hold}} \leftarrow 0$, reset the controller state
\ENDIF
\STATE output $a = a_{\mathrm{hold}}$
\end{algorithmic}
\end{algorithm}

The \ac{qp} is thus solved in the first call after balancing is enabled and
then in every $n_{\mathrm{MPC}}$-th call; between solutions, the last
acceleration is held, which corresponds to the zero-order hold assumed in the
discretisation. When balancing is disabled, the controller state and the
output are reset, and the integrator \eqref{eq:integrator} is reset on the
rising edge of the balancing signal, so that the \ac{mpc} starts from zero
speed and from the current estimate every time the pendulum is captured.

\subsection{Simulation Model}
\label{subsec:simmodel}

The hardware blocks of the Lucas-Nülle toolbox require the controller and its
licence dongle. A simulation copy of the course model,
\emph{A62\_MPC\_Sim}, was therefore created in which the hardware blocks are
replaced by a plant subsystem, while the Control-Routine is identical to the
one used on the hardware. The plant subsystem implements the equations of the
\emph{Position and Angle Control System} block of the Lucas-Nülle simulation
library~\cite{ln_manual} (Figure~\ref{fig:plantsim}):
\begin{itemize}
  \item the drive: the voltage $y_1$ is converted into a speed set-point
        ($0.5/10$~m/s per V), limited to $\pm 0.5$~m/s, and tracked by a
        discrete PI speed loop ($K_p = 100$, $K_i = 30$, output limited to
        $\pm 10~\mathrm{m/s^2}$, back-calculation anti-windup), whose output is
        the cart acceleration;
  \item the pendulum: the nonlinear equation \eqref{eq:pend_nl}, extended by a
        viscous friction term,
        \begin{equation}
          \ddot{\theta} = \frac{-ml\,\ddot{x}\cos\theta - mgl\sin\theta
                                - I f_r\dot{\theta}}{I + ml^2},
          \qquad f_r = 0.3,
          \label{eq:plant_sim}
        \end{equation}
        which is not contained in the prediction model;
  \item the sensors: quantisation of the position to $0.054/16000$~m and of
        the angle to $2\pi/65535$~rad, and the conversion of $\theta$ into the
        wrapped angle $\varphi = \operatorname{atan2}(\sin(\theta-\pi),
        \cos(\theta-\pi))$, which is zero upright.
\end{itemize}
The plant is integrated with the forward Euler method at a fixed step of
$T_{\mathrm{plant}} = 125~\mu\mathrm{s}$, independently of the rate of the
Control-Routine, so that changes of the controller rate do not change the
accuracy of the plant simulation. The model therefore includes the following
differences from the prediction model: the nonlinearity of the pendulum, the
friction term, the dynamics and limits of the drive speed loop, sensor
quantisation, and the discretisation of the controller.

\placeholderfig{0.9}{a62_mpc_sim_plant.png}
  {Plant subsystem of the simulation model A62\_MPC\_Sim, replacing the hardware blocks of the CO3620-2G.}
  {fig:plantsim}
  {Screenshot of the subsystem ``CO3620-2G Plant'' in A62\_MPC\_Sim.slx (drive speed loop PI, integrators, thetaddot Fcn, quantizers, CAN angle). Please provide.}

\subsection{Real-Time Considerations}
\label{subsec:realtime}

At 8~kHz, every call of the Control-Routine must complete within 125~$\mu$s,
including the call in which the \ac{qp} is solved. The solution time of the
\ac{qp} on the TMS320F28335 has not been measured yet. To provide a larger
time budget, the Control-Routine frequency of the hardware block is to be
reduced from 8~kHz to 1~kHz, which gives 1~ms per call; with $T_s = 10$~ms,
$n_{\mathrm{MPC}}$ becomes 10. Because the observer and the filters of the
course model are discrete integrators that inherit the sample time of the
Control-Routine, they adapt to the new rate without modification. The effect
of this reduction on the observer, the swing-up, and the closed loop was
examined in the simulation model; the results are reported in the Results
chapter. Should the
processor load still approach its limit, the control horizon can be reduced to
$N_c = 3$ or the Control-Routine frequency to 500~Hz.

%--------------------------------------------------------------------------
\section{Evaluation Procedure}
\label{sec:evaluation}
%--------------------------------------------------------------------------

\subsection{Simulation Scenarios}
\label{subsec:scenarios}

The controller is evaluated in two stages. First, the \ac{mpc} is closed
directly around the linear discrete-time model with the true state fed to the
controller, in order to verify the formulation, the constraint handling, and
the feasibility of the \ac{qp}. Second, the complete Control-Routine is
simulated in \emph{A62\_MPC\_Sim} with the nonlinear plant, the observer, and
the mode logic. The scenarios are listed in Table~\ref{tab:scenarios}.

\begin{table}[ht]
\centering
\caption{Simulation scenarios.}
\label{tab:scenarios}
\begin{tabular}{p{0.25\textwidth}p{0.25\textwidth}p{0.4\textwidth}}
\toprule
Scenario & Model & Purpose \\
\midrule
100~mm position step & Linear, state feedback & Nominal tracking performance; tuning (Table~\ref{tab:candidates}) \\
Push, $\dot{\varphi}(0) = 1.5$~rad/s & Linear, state feedback & Disturbance rejection; activation of the acceleration and speed constraints \\
140~mm position step & Linear, state feedback & Large set-point close to the travel limit \\
Swing-up and 100~mm step & A62\_MPC\_Sim (nonlinear) & Capture by the mode logic, bumpless start, tracking with observer and drive model \\
Control-Routine at 1~kHz & A62\_MPC\_Sim (nonlinear) & Effect of the reduced rate planned for the hardware \\
\bottomrule
\end{tabular}
\end{table}

\subsection{Performance Metrics}
\label{subsec:metrics}

The following quantities are evaluated for each scenario:
the rise time (10--90\,\%) and the 2\,\% settling time of the cart position;
the overshoot and the undershoot of the cart position; the maximum absolute
pendulum angle $\max|\varphi|$; the maximum absolute cart speed and
acceleration, $\max|\dot{x}|$ and $\max|a|$, which are compared with the limits
R3--R5; the root mean square of the acceleration as a measure of control
effort; and, for the \ac{mpc}, the maximum number of \ac{qp} iterations and
whether a feasible solution was found at every sample. On the hardware, the
processor load reported by the course model is recorded in addition.

\subsection{Planned Hardware Experiments}
\label{subsec:hwplan}

The hardware experiments have not yet been carried out. They are planned in
the following order, so that the effect of each change can be identified:
\begin{enumerate}
  \item The course model is saved under a new name, the Control-Routine
        frequency is set to 1~kHz, and the \ac{mpc} block, the
        \texttt{useMPC} constant, and the switch are copied from
        \emph{A62\_MPC\_Sim} without changing any Lucas-Nülle block.
  \item The model is run with the course \ac{lqr} (\texttt{useMPC}$\,=0$) to
        verify the reduced rate on the hardware with a known controller, and
        the processor load is recorded.
  \item The \ac{mpc} is activated (\texttt{useMPC}$\,=1$), the processor load
        is recorded, and the scenarios of Table~\ref{tab:scenarios} that can be
        reproduced on the hardware are repeated.
\end{enumerate}

\subsection{Planned Comparison with LQR and PID Control}
\label{subsec:benchmark}

The thesis will compare the \ac{mpc} with an \ac{lqr} and a \ac{pid}-based
controller on the same plant. The course \ac{lqr}
(Section~\ref{subsec:lqr}) is available in the same model through the
\texttt{useMPC} switch, uses the same observer, and therefore allows a
comparison under identical conditions. The \ac{pid}-based controller and the
details of the comparison will be defined in a later stage of the work. The
comparison will use the scenarios and metrics of
Sections~\ref{subsec:scenarios} and~\ref{subsec:metrics}.

%--------------------------------------------------------------------------
\section{Summary}
\label{sec:methodsummary}
%--------------------------------------------------------------------------

A linear, constrained \ac{mpc} with a sample time of 10~ms, a prediction
horizon of 50 steps, a control horizon of 5 steps, and a \ac{dare}-based
terminal cost was designed for the CO3620-2G inverted pendulum. The controller
uses the cart acceleration as input, enforces the acceleration limit as a hard
constraint and the speed and travel limits as soft constraints, and obtains
its state from the existing Luenberger observer. It replaces only the
balancing controller of the course model and can be switched back to the
course \ac{lqr} with a single parameter. The controller has been verified in
simulation with a nonlinear plant model; its execution on the hardware and the
comparison with \ac{lqr} and \ac{pid} control are the next steps of the
work.
